\section{Where Computation Happens: Offloading, Remote Rendering and Latency}\label{sec:compute}

A headset has a few watts of power and a few milliseconds of slack. This section asks which stages of
the sense--render--display pipeline should run on the device, at the edge or in the cloud, and how the
delay that each split adds is budgeted, measured and hidden. The function holds \fnComputeN{} papers,
\fnComputeSh{} of the corpus, and its share has held steady (\fnComputeEarly{} in 2017--2019,
\fnComputeLate{} in 2022--2024). It joins four cells: offloading and remote rendering
(\nxrNProbOffloadingRemoteRendering{} papers), latency (\nxrNProbLatency{}, growing
\nxrCagrProbLatency{} a year against \nxrCagrTotal{} for the corpus), rendered frames
(\nxrNContRenderedFrames{}, whose share rose from \nxrShContRenderedFramesEarly{} to
\nxrShContRenderedFramesLate{}) and AR objects and overlays (\nxrNContArObjectsOverlays{}, whose share
fell from \nxrShContArObjectsOverlaysEarly{} to \nxrShContArObjectsOverlaysLate{}). The evidence is
uneven. Offloading papers divide between real testbeds (\nxrRealOffloadingRemoteRendering{}) and
simulation (\nxrSimOffloadingRemoteRendering{}), and only \nxrUsersOffloadingRemoteRendering{} involve
users. Latency papers are more empirical (\nxrRealLatency{} real testbed, \nxrUsersLatency{} with
users). Radio scheduling and transport are treated in \cref{sec:provision}, and QoE models in
\cref{sec:assure}.

\begin{figure}[t]
\centering
\begin{tikzpicture}[x=1cm, y=1cm, font=\sffamily\scriptsize,
  st/.style={draw=compute!70!black, fill=compute!10, rounded corners=1.5pt, line width=0.5pt,
             minimum width=2.05cm, minimum height=0.62cm, align=center, inner sep=1.5pt},
  est/.style={st, minimum width=2.3cm},
  loc/.style={-{Stealth[length=1.8mm]}, line width=0.5pt, draw=black!75},
  net/.style={-{Stealth[length=2.2mm]}, line width=1.1pt, draw=compute!75!black},
  alt/.style={-{Stealth[length=1.8mm]}, line width=0.5pt, dashed, draw=black!55},
  lab/.style={font=\sffamily\scriptsize, text=black, inner sep=1.5pt},
  lane/.style={font=\sffamily\scriptsize\bfseries, anchor=west, align=left, inner sep=0pt, text=black},
  hdr/.style={font=\sffamily\footnotesize\bfseries, text=black}]
  % ---- tier columns
  \fill[black!3]   (0,-5.65) rectangle (6.85,0.4);
  \fill[black!8]   (6.85,-5.65) rectangle (9.55,0.4);
  \fill[compute!7] (9.55,-5.65) rectangle (12.6,0.4);
  \fill[black!3]   (12.6,-5.65) rectangle (16.3,0.4);
  \node[hdr] at (3.42,0.12)   {Device (headset, phone)};
  \node[hdr] at (8.2,0.12)    {Access network};
  \node[hdr] at (11.075,0.12) {Edge server};
  \node[hdr] at (14.45,0.12)  {Cloud};
  \node[lab, text=black!70] at (8.2,-0.2) {Wi-Fi, 60~GHz, 4G/5G};
  \foreach \y in {-2.15,-3.9} \draw[black!30, line width=0.4pt] (0,\y) -- (12.6,\y);
  % ---- lane 1: full remote rendering
  \node[lane] at (0.12,-0.85) {(1) Full remote\\\phantom{(1) }rendering};
  \node[st]  (s1) at (5.78,-0.85) {Sense pose};
  \node[st]  (d1) at (5.78,-1.70) {Decode\\\textcolor{black!70}{14--16~ms}};
  \node[st]  (w1) at (3.48,-1.70) {Reproject,\\time-warp};
  \node[st]  (p1) at (1.18,-1.70) {Display};
  \node[est] (r1) at (11.075,-0.85) {Render full frame};
  \node[est] (e1) at (11.075,-1.70) {Encode\\\textcolor{black!70}{9--10~ms}};
  \draw[net] (s1) -- node[lab, above] {pose} (r1);
  \draw[loc] (r1) -- (e1);
  \draw[net] (e1) -- node[lab, below] {encoded frames} (d1);
  \draw[loc] (d1) -- (w1); \draw[loc] (w1) -- (p1);
  % ---- lane 2: split rendering
  \node[lane] at (0.12,-2.60) {(2) Split\\\phantom{(2) }rendering};
  \node[st]  (s2) at (5.78,-2.60) {Sense pose};
  \node[st]  (n2) at (3.48,-2.60) {Render near\\objects};
  \node[st]  (d2) at (5.78,-3.45) {Decode};
  \node[st]  (w2) at (3.48,-3.45) {Compose,\\warp};
  \node[st]  (p2) at (1.18,-3.45) {Display};
  \node[est] (r2) at (11.075,-2.60) {Render far scene};
  \node[est] (e2) at (11.075,-3.45) {Encode layer};
  \draw[net] (s2) -- node[lab, above] {pose} (r2);
  \draw[loc] (s2) -- (n2); \draw[loc] (n2) -- (w2);
  \draw[loc] (r2) -- (e2);
  \draw[net] (e2) -- node[lab, below] {far-scene layer} (d2);
  \draw[loc] (d2) -- (w2); \draw[loc] (w2) -- (p2);
  % ---- lane 3: AR detection offload
  \node[lane] at (0.12,-4.35) {(3) AR detection\\\phantom{(3) }offload};
  \node[st]  (s3) at (5.78,-4.35) {Capture camera\\frame};
  \node[st]  (t3) at (5.78,-5.20) {Track locally};
  \node[st]  (o3) at (3.48,-5.20) {Render overlays};
  \node[st]  (p3) at (1.18,-5.20) {Display};
  \node[est, minimum height=1.47cm] (x3) at (11.075,-4.775) {Detect,\\recognise\\(GPU)};
  \draw[net] (s3) -- node[lab, above] {camera frames} (s3 -| x3.west);
  \draw[net] (t3 -| x3.west) -- node[lab, below] {labels, object poses} (t3);
  \draw[loc] (s3) -- (t3); \draw[loc] (t3) -- (o3); \draw[loc] (o3) -- (p3);
  % ---- cloud
  \node[draw=black!45, fill=white, rounded corners=1.5pt, line width=0.5pt, text width=3.05cm,
        minimum height=4.6cm, align=center, inner sep=3pt] (cl) at (14.45,-2.82)
        {Any edge function can\\run here instead\\[8pt]
         Round-trip time 49.3~ms,\\against $\approx$1~ms to a fog\\node or cloudlet\\[8pt]
         Simulated AR latency\\12.85\% above the edge};
  \draw[alt] (r1.east) -- (r1.east -| cl.west);
  \draw[alt] (r2.east) -- (r2.east -| cl.west);
  \draw[alt] (x3.east) -- (x3.east -| cl.west);
  % ---- MTP brace and legend
  \draw[decorate, decoration={brace, mirror, amplitude=4pt}, draw=black!60, line width=0.5pt]
    (0.05,-5.75) -- (12.55,-5.75)
    node[midway, below=5pt, lab, align=center]
    {Motion-to-photon (MTP) loop of lane (1): budget 15--20~ms, of which $\approx$14~ms for computing and network;\\
     achieved: 10--14~ms end to end over 60~GHz with parallel encoding and display-synchronised rendering};
  \begin{scope}[shift={(12.95,-5.92)}]
    \draw[net] (0,0) -- (0.6,0);         \node[lab, anchor=west] at (0.68,0)     {crosses the network};
    \draw[loc] (0,-0.3) -- (0.6,-0.3);   \node[lab, anchor=west] at (0.68,-0.3)  {stays in one tier};
    \draw[alt] (0,-0.6) -- (0.6,-0.6);   \node[lab, anchor=west] at (0.68,-0.6)  {alternative host};
  \end{scope}
\end{tikzpicture}
\caption{Offloading architecture: where each stage runs in full remote rendering, split rendering and
AR detection offload. Each split is the point at which the pipeline crosses the access network;
display, and in AR also tracking, stay on the device in every design. Values: MTP budget \cite{W2784711283}, 60~GHz end-to-end latency
\cite{W2838572662}, encode and phone decode times \cite{W2940859338}, round-trip times
\cite{W2805833065}, simulated edge advantage \cite{W2894645340}.}
\label{fig:compute-pipeline}
\end{figure}

\subsection{Choosing the split}

Proximity is the one placement result that holds across methods. In a shipyard deployment, a fog
gateway and a cloudlet both answered in about 1~ms, against 49.3~ms for a remote cloud; at 240
clients the fog gateway was four times slower than the cloudlet, and the Wi-Fi hop was excluded by
design \cite{W2805833065}. A simulation found that edge execution cut AR latency by 41.44\% against
local and 12.85\% against cloud execution \cite{W2894645340}. What leaves the device varies
(\cref{fig:compute-pipeline}). VR systems render remotely and stream frames. Split designs render near
or interactive objects on the device and the rest remotely; one used 20--30\% of the bandwidth of full
remote rendering and tolerated delays of up to 200~ms \cite{W4399444696,W4389296467,W3080913128}. AR
systems offload perception and render locally. Neural scenes reopen the question: current systems
stream pixels \cite{W4415672032}, dynamic Gaussians \cite{W7106240991} or intermediate level-of-detail
results \cite{W7117245976}, each against its own baseline.

Shared edges add contention, and here the evidence thins. Once many users offload, queueing at the
edge server dominates \cite{W4213053765,W4399320016}, yet schedulers for concurrent rendering sessions
on one GPU exist only in simulation \cite{W3214338661,W4285254743}. The largest strand formulates
offloading as joint computation--communication optimisation, solved by deep reinforcement learning
\cite{W3120265095} or Lyapunov control \cite{W4392567309}. More than nine in ten learning-based
controllers are evaluated only in simulation, and their inputs are assumed: a 2018 YOLO accuracy fit
\cite{W2790412025} reappears in later formulations \cite{W4386702674,W4388068949}, and VR streams are
modelled at a constant 25.32~Mbps \cite{W2596760110}.

Device energy, the main reason to offload, is contested. Models promise savings of 63\%
\cite{W2608841173} and 85\% \cite{W4388068949}. Measurements disagree. Edge offloading halved
smartphone energy for AR on average \cite{W2983318720}, and an edge VR testbed lowered client power
\cite{W4296339041}. A Quest Pro, however, drained its battery 15\% faster with wireless offloading
than with local rendering \cite{W4362489903}, and Air Link on a Quest~2 also cost more energy
\cite{W4387913063}. DeepDecision measured about 2900~mW for offloading regardless of bitrate and frame
rate \cite{W2792220137}, against models that scale energy with transmitted data. Radio tail, decoder
and display power, omitted from most models, are the likely cause.

\subsection{Remote rendering: frames, codecs and pacing}

Compression is mandatory, and its cost is set by pipelining and timing more than by link rate. Raw
stereo frames at 2160$\times$1200 and 90~Hz need 5.6~Gbps. Cutting the Cord split each eye's frame
across four parallel H.264 encoders and aligned rendering with the client's display refresh; over
802.11ad it reached 10--14~ms end to end, against a sequential median above 26~ms, and cut missed
refreshes from 5.3--14.3\% to 0.1--0.2\% \cite{W2838572662}. Phones are tighter. Furion measured
9--10~ms to encode and 14--16~ms to decode H.264 frames, and 27\% of the phone's CPU spent on packet
processing at 400~Mbps \cite{W2940859338}. Removing memory copies brought OpenUVR to 14.32~ms
\cite{W3124215200}.

The renderer knows what the encoder must otherwise estimate, and using it pays. Object-aware
quantisation with gaze saved 40--47\% bandwidth in a subjective test \cite{W3174397927}. Foveated
shading alone did not reduce bitrate, while foveated post-processing or encoding did
\cite{W3186601312}. Other systems change what is streamed: colour plus depth for client-side warping
\cite{W2751899584}, object-space shading atlases \cite{W2903169695}, one eye rendered remotely and the
other reprojected with 36\% lower MTP \cite{W4391557896}, and pixel reuse across frames with 95\% less
frame data \cite{W2928088097}. Each uses its own baseline, so they cannot be ranked (\cref{sec:encode}).
Two assumptions remain untested. No corpus paper measures AV1 or VVC latency for rendered frames;
the only mention is a position paper \cite{W4391267710}. Frame sizes are bursty even under constant-bitrate coding \cite{W4379468406}, and
on one headset the 3GPP traffic model underestimated effective bitrate by 25\% \cite{W4399340506}.

\subsection{AR offloading: stale results and the uplink}

AR offloading succeeds only if the device hides the staleness of remote results. The replicated
design runs a light tracker on the device between remote detections. In one system, without
tracking, 16.7~ms of offload latency raised false detections from 0 to 31.56\%; with tracking and
region-of-interest encoding, offload latency fell from 34.56 to 17.23~ms on 2.4~GHz Wi-Fi
\cite{W2965289829}. CloudAR held 30 frames per second with at least 97\% recognition accuracy
\cite{W2765872303}, Jaguar reached about 33~ms end to end \cite{W2896416526}, and the pattern holds on
HoloLens \cite{W4282961825}. Once the server has a GPU, the network dominates: sending frames as H.264
video rather than images delivered about ten times more frames at 100~kbps \cite{W2792220137}.

The uplink is the bottleneck. In multi-user AR over LTE, median end-to-end latency was 12.5~s, the
radio access network contributed 31.2\%, and uplink bursts triggered TCP slow starts
\cite{W3047439371}. A 5G mmWave measurement found the uplink insufficient for baseline edge detection
\cite{W4390905995}, while vision papers assume 5G solves the problem \cite{W3133913697}. Result sharing among co-located users cut
latency by up to 130$\times$ \cite{W4213053765}, but the overlap it relies on has never been measured. Targets are
contested: below 20~ms \cite{W2609285553}, 50~ms \cite{W3133913697} and 100~ms
\cite{W2898344194,W4282961825} are all quoted. The right one depends on whether overlays are
world-locked and tracked locally, and no study maps network delay to registration error.

\subsection{Budgeting, measuring and hiding delay}

Latency budgets are inherited rather than derived. The common figure is 15--20~ms MTP, leaving about
14~ms for computing and communication \cite{W2784711283}; others assume 13~ms \cite{W2561715836} or
5~ms \cite{W3048503162}. Users are more tolerant once reprojection is active. CloudVR users did not
notice 20~ms of added network latency \cite{W2896076858}, and ALVR users hardly noticed round-trip
times up to 90~ms \cite{W3095344018}, whereas CloudXR quality fell from 75~ms \cite{W4402156718}. In a
crane simulator, display delays of 0--30~ms mattered while controller delay mattered only at 800~ms
\cite{W3086540722}. A budget therefore depends on which loop is meant and how delay is hidden, and
offloading papers rarely state either. Measurement is cheap but rare: a photosensor rig
\cite{W2890615649} and a phone-camera method with under 10~ms uncertainty \cite{W3012329178} measure
MTP directly, yet fewer than one in twenty latency papers report jitter or tail latency.

Latency hiding is the established success. Post-render warping removed up to 80\% of the aiming
penalty of 80~ms added latency \cite{W3089399051}. Rendering beyond the field of view, sized to the
current latency, reduced delivery latency by about 22\% in a simulation with real user data
\cite{W3161249521}. Motion predictors trained on more than 840,000 samples let the edge pre-render
frames \cite{W3094082851}. Reprojection has moved from rotation to translation \cite{W4408848620}, and
selective time-warp leaves head-locked layers untouched \cite{W2908115098}. Stalls of up to four frames
at 90~Hz went unnoticed, and reprojection helped from eight frames upward \cite{W4301184864}. Each
method, however, is tested on its own motion data and with fixed delays, so the gains cannot be
compared.

\begin{openproblems}
\item \textbf{Learning-based offloading controllers are untested on real traces.} More than nine in
  ten are evaluated in simulation with synthetic task sizes and channels. An open benchmark replaying
  measured render times, head motion and 5G/Wi-Fi traces under 3GPP XR traffic would show which gains
  survive.
\item \textbf{Device energy savings are modelled, not measured.} One Quest Pro drained 15\% faster
  with offloading \cite{W4362489903}; one edge testbed lowered power \cite{W4296339041}. Component-level
  measurements of local and remote rendering across headsets, codecs and links would settle the
  direction.
\item \textbf{3GPP XR traffic models go unused and unvalidated.} No offloading paper evaluates against
  them, and the one empirical check found a 25\% bitrate underestimate \cite{W4399340506}. Commercial
  split-rendering traces fitted against the models are needed.
\item \textbf{Rendered-frame streams lack error resilience.} Users rank loss above delay
  \cite{W3095344018,W4404916373}, yet forward error correction, intra refresh and concealment for
  remote rendering barely appear in the corpus. A loss-trace study with user ratings would close this.
\item \textbf{No benchmark couples motion traces, network traces and MTP ground truth.} Warping and
  prediction methods each use their own data and fixed delays. A shared dataset with photodiode MTP
  would let them be ranked.
\item \textbf{The AR camera uplink is unstudied at multi-user scale.} Real-network evidence is
  single-user and pessimistic \cite{W4390905995,W3047439371}. A 5G testbed scheduling several AR
  uplinks in one cell would test whether edge AR scales.
\end{openproblems}
